\documentclass[sigconf]{acmart}

% Copyright information
\setcopyright{acmcopyright}
\acmYear{2025}
\acmConference[CONF'25]{Conference Name}{Month Day-Day, 2025}{City, Country}
\acmBooktitle{Conference Name (CONF'25), Month Day-Day, 2025, City, Country}
\acmDOI{xx.xxxx/xxxxxxx.xxxxxx}

% Packages
\usepackage{booktabs}
\usepackage{graphicx}
\usepackage{amsmath}
\usepackage{algorithm}
\usepackage{algorithmic}
\usepackage{xcolor}
\usepackage{hyperref}

% Document starts
\begin{document}

% Title and subtitle
\title{Title of Your Research Paper}
\subtitle{Optional Subtitle}

% Author information
\author{Author Name}
\affiliation{
    \institution{Institution Name}
    \department{Department Name}
    \city{City}
    \country{Country}
}
\email{author@institution.edu}

\author{Second Author}
\affiliation{
    \institution{Institution Name}
    \department{Department Name}
    \city{City}
    \country{Country}
}
\email{second@institution.edu}

% Abstract
\begin{abstract}
This is the abstract of your paper. It should be 150-200 words and summarize the 
purpose, methods, results, and conclusions of your research. The abstract should 
be self-contained and readable without the full paper. Include specific metrics 
and quantitative results where possible.
\end{abstract}

% CCS Concepts
\begin{CCSXML}
<ccs2012>
   <concept>
       <concept_id>10010147.10010178</concept_id>
       <concept_desc>Computing methodologies</concept_desc>
       <concept_significance>500</concept_significance>
       </concept>
   <concept>
       <concept_id>10010147.10010178.10010224</concept_id>
       <concept_desc>Computing methodologies~Computer vision</concept_desc>
       <concept_significance>500</concept_significance>
       </concept>
 </ccs2012>
\end{CCSXML}

\ccsdesc[500]{Computing methodologies}
\ccsdesc[500]{Computing methodologies~Computer vision}

% Keywords
\keywords{keyword1, keyword2, keyword3, keyword4}

% Make title
\maketitle

% Section 1: Introduction
\section{Introduction}
\label{sec:introduction}

Start with a compelling real-world problem or motivation. Build from general to 
specific context. Clearly state the research gap and your contributions.

The main contributions of this paper are:
\begin{itemize}
    \item Contribution 1: Description of your first contribution.
    \item Contribution 2: Description of your second contribution.
    \item Contribution 3: Description of your third contribution.
\end{itemize}

The remainder of this paper is organized as follows. Section~\ref{sec:related} 
reviews related work. Section~\ref{sec:method} describes the proposed method. 
Section~\ref{sec:experiments} presents experimental results. 
Section~\ref{sec:conclusion} concludes the paper.

% Section 2: Related Work
\section{Related Work}
\label{sec:related}

Review relevant literature and position your work. Group related work by theme 
or approach. Compare explicitly: "Unlike [X] which focuses on Y, our approach 
addresses Z..."

% Section 3: Methodology
\section{Methodology}
\label{sec:method}

Describe your proposed method in detail. Include architecture diagrams, 
algorithms, and design rationale.

\subsection{Architecture Overview}

Provide a high-level description of your system or method.

\subsection{Key Components}

Describe each component in detail.

\subsection{Algorithm}

Present key algorithms using pseudocode:

\begin{algorithm}
\caption{Algorithm Name}
\label{alg:algorithm}
\begin{algorithmic}[1]
\REQUIRE Input parameters
\ENSURE Output results
\STATE Step 1
\STATE Step 2
\STATE Step 3
\RETURN result
\end{algorithmic}
\end{algorithm}

% Section 4: Experiments
\section{Experiments}
\label{sec:experiments}

\subsection{Experimental Setup}

Describe datasets, evaluation metrics, baselines, and implementation details.

\subsection{Results}

Present quantitative results in tables:

\begin{table}[htbp]
\caption{Performance Comparison}
\label{tab:results}
\centering
\begin{tabular}{@{}llll@{}}
\toprule
Method & Metric 1 & Metric 2 & Metric 3 \\
\midrule
Baseline & 0.80 & 0.75 & 0.82 \\
Proposed & \textbf{0.91} & \textbf{0.89} & \textbf{0.92} \\
\bottomrule
\end{tabular}
\end{table}

\subsection{Analysis}

Provide ablation studies, error analysis, and qualitative examples.

% Section 5: Conclusion
\section{Conclusion}
\label{sec:conclusion}

Summarize your contributions and their significance. Discuss limitations and 
future work.

% Acknowledgments (optional)
\begin{acks}
Thank funding sources, collaborators, etc.
\end{acks}

% References
\bibliographystyle{ACM-Reference-Format}
% \bibliography{references}

% Example manual references (or use .bib file)
\begin{thebibliography}{1}

\bibitem{example1}
A. Author and B. Author. Year. Title of paper. In \emph{Proceedings of 
Conference Name} (CONF'YY). ACM, New York, NY, USA, 123--456. 
\url{https://doi.org/xx.xxxx/xxxxx}

\bibitem{example2}
C. Author. Year. Title of journal article. \emph{Journal Name} X, Y (Month 
Year), 123--456. \url{https://doi.org/xx.xxxx/xxxxx}

\end{thebibliography}

\end{document}
